\section{Анализ работ}
\subsection{Смежные исследования и исходная проблема}
В студенческой работе С.~Аскари Хеммат рассматривается подготовка синтетических сцен с мебелью в Unreal Engine и обучение моделей на этих данных~\cite{askari}. Она связана с настоящей темой способом получения изображений, но приведённые в ней показатели детекции не являются оценкой UDA на дорожных сценах. Л.~Б.~Комендала рассматривает сегментацию дорожных изображений CamVid~\cite{komendala}; эта работа близка по прикладной задаче, однако не задаёт перенос GTA5~$\rightarrow$~Cityscapes. В исследовании И.~И.~Курочкина и соавторов нейросети применяются для классификации микроскопических изображений криогелей~\cite{kurochkin}. Это пример анализа изображений в другой предметной области; его результаты не переносятся на метрики семантической сегментации.

\subsection{Основные направления адаптации}
Статистические критерии сравнивают распределения признаков через ядерные средние или ковариации~\cite{mmd,coral}. Оптимальный транспорт строит соответствия между наблюдениями, минимизируя стоимость переноса массы~\cite{sinkhorn}. Состязательные методы обучают модель скрывать различие доменов от дополнительного классификатора. В AdaptSegNet согласование проводится на уровне пространственно организованных предсказаний сегментатора~\cite{tsai}.

Другая группа подходов использует структуру ответов на целевых изображениях. Энтропийный критерий поощряет определённые предсказания~\cite{advent}. Самообучение превращает часть предсказаний в обучающие метки. Согласованность требует устойчивости ответов к допустимым преобразованиям; в MIC ученик предсказывает классы по частично закрытому изображению, используя псевдометки учителя~\cite{mic}. Контрастивные методы, например SePiCo, добавляют семантические ограничения на расположение признаков~\cite{sepico}.

Направления пересекаются. В DACS псевдометки используются вместе со смешиванием исходных и целевых изображений и схемой «учитель --- ученик». DAFormer дополняет самообучение выборкой редких классов, регуляризацией признаков и разогревом скорости обучения. Поэтому название алгоритма обозначает целую реализацию, а не единственный математический критерий. В работе рассматриваются шесть перечисленных направлений; преобразование внешнего вида изображений не выделяется в самостоятельный предмет сравнения.

\subsection{Что показывают опубликованные оценки}
Для оценки всей схемы нужна модель, обученная только на исходном домене (source only). В AdaptSegNet многоуровневое согласование повышает mIoU с 36,6 до 42,4, то есть на 5,8 процентного пункта (п.~п.)~\cite[табл.~1]{tsai}. В DACS соответствующая пара равна 32,85 и 52,14, прирост составляет 19,29~п.~п.~\cite[табл.~1]{dacs}. Обе оценки относятся к GTA5~$\rightarrow$~Cityscapes и 19 классам, но условия обучения различаются. Разница приростов не доказывает преимущество самообучения над состязательным критерием в одинаковой модели.

Внутреннее сравнение DACS (табл.~\ref{tab:dacs}) показывает риск ухудшения при псевдоразметке. Naive Mixing смешивает целевые изображения между собой, DACS --- исходные и целевые. Различие результатов характеризует эти реализации, а не направление самообучения в целом.

\begin{table}[htbp]
\caption{Опубликованное сравнение вариантов DACS на GTA5~$\rightarrow$~Cityscapes~\cite[табл.~3]{dacs}}
\label{tab:dacs}
\centering
\setlength{\tabcolsep}{3pt}
\begin{tabularx}{\columnwidth}{@{}Xrr@{}}
\toprule
Вариант & mIoU, \% & $\Delta$, п.~п.\\
\midrule
Только исходный домен & 32,85 & ---\\
Псевдоразметка без смешивания & 22,97 & $-9{,}88$\\
Naive Mixing & 35,08 & $+2{,}23$\\
DACS & 52,14 & $+19{,}29$\\
\bottomrule
\end{tabularx}
\par\vspace{3pt}
{\footnotesize $\Delta$ рассчитана относительно 32,85. Это результаты авторов DACS, не собственные эксперименты.}
\end{table}

Эффект отдельного компонента оценивают относительно уже адаптирующейся системы. Добавление MIC к DAFormer изменяет mIoU с 68,3 до 70,6, то есть на 2,3~п.~п.~\cite[табл.~7]{mic}. Это дополнительный эффект MIC; исходные 68,3~mIoU~\cite{daformer} характеризуют сочетание архитектуры и правил обучения. Сопоставимого численного примера для CORAL в выбранном сценарии не приводится.
